% annexes/manuel.tex
% Manuel utilisateur - Analyse du reseau WETH Polygon

\section{Manuel utilisateur}
\label{ann:manuel}

Ce manuel d\'ecrit les pr\'erequis, la proc\'edure d'installation, les commandes d'utilisation et la structure des fichiers du projet.

% ---- Acces en ligne ----

\subsection{Acc\`es en ligne au prototype interactif}

Pour explorer les analyses sans rien installer, le prototype web est d\'eploy\'e \`a l'adresse~:

\begin{center}
\large\url{https://crypto-graph-explorer.vercel.app}
\end{center}

Cette version contient les six questions de recherche et les cinq analyses compl\'ementaires de l'annexe~\S\ref{sec:complementaires}, avec visualisation D3.js du r\'eseau, panneau de statistiques en temps r\'eel et simulations interactives (suppression de n\oe{}uds, propagation SI). Les chiffres exacts du rapport (Gini, $\sigma$, modularit\'e $Q$) restent obtenus uniquement par l'ex\'ecution Python locale d\'ecrite ci-apr\`es, le prototype op\'erant sur des donn\'ees pr\'e-calcul\'ees et des m\'etriques approch\'ees pour rester r\'eactif dans un navigateur.

% ---- Prerequis ----

\subsection{Pr\'erequis}

Les \'el\'ements suivants doivent \^etre install\'es sur la machine avant d'utiliser le projet~:

\begin{itemize}
  \item \textbf{Python 3.10 ou sup\'erieur}~: test\'e avec les versions 3.10, 3.11 et 3.12. La version 3.13 est compatible mais certaines biblioth\`eques C peuvent n\'ecessiter une compilation manuelle.
  \item \textbf{pip}~: gestionnaire de paquets Python (inclus avec Python 3.10+).
  \item \textbf{TeX Live 2024} (ou MiKTeX)~: distribution LaTeX compl\`ete incluant pdfLaTeX et Biber pour la compilation du rapport.
  \item \textbf{Git}~: pour le clonage du d\'ep\^ot et le contr\^ole de version.
  \item \textbf{Espace disque}~: au minimum 100~Mo (CSV de 74~Mo + figures g\'en\'er\'ees + rapport PDF).
\end{itemize}

% ---- Installation ----

\subsection{Installation}

Cloner le d\'ep\^ot et installer les d\'ependances Python~:

\begin{lstlisting}[language={}, numbers=none, caption={Clonage et installation}]
$ git clone <repository-url>
$ cd polygon-weth-crash-analysis
$ pip install -r requirements.txt
\end{lstlisting}

\paragraph{Note importante~:} le fichier de donn\'ees \texttt{polygon\_weth\_crash\_2024-08-05.csv} (74~Mo, 401\,709 transactions) doit \^etre plac\'e \`a la racine du projet. Ce fichier n'est pas inclus dans le d\'ep\^ot Git en raison de sa taille. Il peut \^etre r\'eg\'en\'er\'e \`a partir de Google BigQuery~\cite{bigquery_public} via la commande \texttt{python main.py} (n\'ecessite un fichier de cl\'e de service \texttt{google\_key.json}).

% ---- Utilisation ----

\subsection{Utilisation}

Le projet s'utilise en trois \'etapes principales~:

\paragraph{1. Extraction des donn\'ees (optionnel)~:} cette \'etape n'est n\'ecessaire que si le fichier CSV n'est pas d\'ej\`a disponible. Elle requiert un compte Google Cloud Platform et un fichier de cl\'e de service.

\begin{lstlisting}[language={}, numbers=none, caption={Extraction des donn\'ees depuis BigQuery}]
$ python main.py
\end{lstlisting}

\paragraph{2. Ex\'ecution des analyses~:} le script orchestrateur lance s\'equentiellement les six analyses et g\'en\`ere les 16 figures dans \texttt{figures/}.

\begin{lstlisting}[language={}, numbers=none, caption={Lancement des analyses}]
$ python run_all.py
\end{lstlisting}

\paragraph{3. Compilation du rapport~:} g\'en\`ere le fichier PDF final contenant le rapport complet avec figures int\'egr\'ees.

\begin{lstlisting}[language={}, numbers=none, caption={Compilation du rapport LaTeX}]
$ cd rapport
$ latexmk -pdf main.tex
\end{lstlisting}

% ---- Structure des fichiers ----

\subsection{Structure des fichiers}

L'arborescence compl\`ete du projet est la suivante~:

\begin{lstlisting}[language={}, numbers=none, caption={Arborescence du projet}]
projet/
+-- polygon_weth_crash_2024-08-05.csv  (donnees, 74 Mo)
+-- run_all.py                         (orchestrateur)
+-- main.py                            (fetching BigQuery)
+-- requirements.txt                   (dependances Python)
+-- src/
|   +-- graph_builder.py               (construction du graphe)
|   +-- style.py                       (style matplotlib)
|   +-- rq1_topologie.py               (RQ1 : topologie)
|   +-- rq2_centralite.py              (RQ2 : centralite)
|   +-- rq3_communautes.py             (RQ3 : communautes)
|   +-- rq4_temporel.py                (RQ4 : temporel)
|   +-- rq5_petitmonde.py              (RQ5 : petit-monde)
|   +-- rq6_flux_ponderes.py           (RQ6 : flux ponderes)
+-- figures/                           (16 PNG, 200 DPI)
+-- rapport/
    +-- main.tex                       (document principal)
    +-- references.bib                 (bibliographie)
    +-- sections/                      (14 fichiers .tex)
    +-- annexes/                       (3 fichiers .tex)
\end{lstlisting}

\paragraph{Description des r\'epertoires~:}
\begin{itemize}
  \item \texttt{src/}~: contient les 8 modules Python. Chaque module \texttt{rq\{1-6\}\_*.py} expose une fonction \texttt{run(df, digraph, graph)} qui ex\'ecute l'analyse correspondante et g\'en\`ere les figures associ\'ees.
  \item \texttt{figures/}~: r\'epertoire de sortie des 12 visualisations PNG \`a 200~DPI. Les figures \`a fond sombre (\texttt{\#0F172A}) repr\'esentent les graphes r\'eseau, tandis que les figures \`a fond clair correspondent aux graphiques analytiques.
  \item \texttt{rapport/}~: contient le document LaTeX principal (\texttt{main.tex}), le fichier de bibliographie (\texttt{references.bib}), les 14 fichiers de sections et les 3 fichiers d'annexes.
\end{itemize}
